\chapter{Village Studies}

\section{Is Caste a Closed Class or an Open Caste?}

\begin{KeyConcept}
\begin{itemize}
\item On 'is caste changing/weakening/disintegrating?' --- the old belief was that \textbf{caste is closed and class is open}; but \textbf{caste is a closed class and class is an open caste}.
\item \textbf{Pitirim Sorokin}: \textbf{no society is absolutely closed or absolutely open --- it is a matter of degree.} Caste does not disappear: \textbf{mobility increases, but caste persists} (today the upper caste can be found in lower occupations and the lower caste in higher occupations --- but occupation changes while the caste remains, and it reappears at marriage).
\end{itemize}
\end{KeyConcept}

\section{Questioning the Self-Sufficiency Doctrine}

\begin{KeyConcept}
\begin{itemize}
\item In the \textbf{1950s}, ethnographers began to \textbf{question the self-sufficiency doctrine} of the Orientalists (Metcalfe, Marx, the Indologists).
\item \textbf{M. N. Srinivas and A. M. Shah} wrote an EPW article, \textbf{'Myth of Self-Sufficiency'}, questioning it through a \textbf{Gujarat case study} (trading; migration between Gujarat, Maharashtra and MP).
\item \textbf{Beteille's} Shripuram study showed the same: it had \textbf{Kannada-speaking, Telugu-speaking and Tamil-speaking Brahmins} --- clear evidence that villagers always had \textbf{connections with outsiders} (Harappan seals found in Mesopotamia; Chola/Pallava trade with the Roman Empire; land grants sold for migration).
\item \textbf{S. C. Dube}: the Indian village has always been \textbf{dependent on towns} for salt, bangles and things of daily use; towns depended on villages for manpower, food-grains and raw materials --- so \textbf{rural and urban India were interdependent and interrelated}.
\item \textbf{Jyotiba Phule}: the \textbf{drain of wealth} (Dadabhai Naoroji's theory, given with respect to two nations) operates \textbf{within India too} --- the same exploitation of rural India by urban India.
\end{itemize}
\end{KeyConcept}

\section{S. C. Dube: A Typology of Villages}

\begin{KeyConcept}
\begin{itemize}
\item S. C. Dube conducted three field studies: the \textbf{Kamars} (Chhattisgarh, then MP), \textbf{Shamirpet} (Telangana, 1955) and \textbf{Ranikhera} (UP). He focused on \textbf{Shamirpet}.
\item Shamirpet was chosen because it is a \textbf{multi-caste and multi-religious village} --- Hindus, Muslims, Christians, Dalit Christians, Dalit Muslims; Telugu, Hindi and Urdu speakers (508 households; 1,437 Hindu upper castes; 680 SC/ST; 340 Muslims; 2,008 Telugu speakers; 340 Urdu speakers). It is 'India in miniature'.
\item Dube observed \textbf{jajmani across religious groups}, not just caste groups: \textbf{Hindus and Muslims are interdependent} for goods and services (e.g. the food of Old Delhi is a symbol of Hindu--Muslim integration) --- this maintained \textbf{integration and consensus} in the village.
\end{itemize}
\end{KeyConcept}

\begin{KeyConcept}
\begin{itemize}
\item Dube's point: \textbf{there is no single 'Indian village'} (a Punjabi village differs from a Himachali, Kashmiri, Keralite or Telangana village). The colonials and social anthropologists developed a \textbf{monolithic image} of the Indian village without doing field work.
\item So instead of theorising the Indian village, Dube developed a \textbf{typology} --- \textbf{six criteria} to differentiate villages: (1) \textbf{area of land and size of population}, (2) \textbf{ethnic and caste composition}, (3) \textbf{pattern of land ownership}, (4) \textbf{structure of authority and power hierarchy} (who is the headman and on what basis), (5) \textbf{proximity to towns} (closer = more contact/assimilation/urbanisation), (6) \textbf{local traditions} (deities, festivals, their functions).
\end{itemize}
\end{KeyConcept}

\begin{KeyConcept}
\begin{itemize}
\item From these six criteria, \textbf{three broad features} of the Indian village emerge: (1) it is a \textbf{territorial, social, political and economic unit}; (2) there is \textbf{inter-ethnic (inter-caste) dependency}; (3) \textbf{multiple villages are integrated with each other} (not self-sufficient).
\item This is a \textbf{structural-functional} model --- Dube saw dependency and integration, but \textbf{not conflict}.
\end{itemize}
\end{KeyConcept}

\section{S. C. Dube: How the Village Community is Organized}

\begin{KeyConcept}
\begin{itemize}
\item The village community is organised in \textbf{two ways}:
\item \textbf{1. In terms of kin and caste} --- the individual's first obligation is to the \textbf{family}, then the \textbf{group of near kin}, then the \textbf{lineage} (vansh --- the descent symbol, e.g. the Yadavs tracing descent to Yadu Maharaj), then \textbf{relatives} (affinal --- by marriage), then the \textbf{sub-caste/caste}, then the \textbf{varna}. (In Beteille's Shripuram: family $\rightarrow$ kin $\rightarrow$ lineage $\rightarrow$ relatives $\rightarrow$ Smartha/Sri Vaishnava sub-caste $\rightarrow$ Brahmin --- activated in different contexts.)
\item \textbf{2. In terms of territorial affiliation} --- the individual's affiliation is to the \textbf{family}, then the \textbf{village}, then the \textbf{intra-caste organization}, then the \textbf{region}, then the \textbf{nation}.
\item Overall, Dube tried to \textbf{develop models} of the Indian village, not a monolithic image.
\end{itemize}
\end{KeyConcept}

\section{Beteille's Weberian Model of Social Change}

\begin{KeyConcept}
\begin{itemize}
\item \textbf{Beteille} (a Weberian) did not develop a generalisation: he developed a \textbf{model of social change} from his Shripuram study --- from \textbf{symmetrical} (caste = class = power) to \textbf{asymmetrical} (class and power disembedded from caste).
\item This model is an \textbf{ideal type}, to be used for \textbf{comparative analysis} --- to see the \textbf{divergences and convergences} between the model and the reality of other villages (every village is unique; the social world is not so patterned that we can generalise from one village to all India).
\end{itemize}
\end{KeyConcept}

\section{Marxist Village Studies: Thorner and Gough}

\begin{KeyConcept}
\begin{itemize}
\item \textbf{From the 1970s} (after the Green Revolution, 1967--77), village studies witnessed the \textbf{advent of Marxist studies} --- which \textbf{rejected the unity/integration principle} of structural functionalism and brought the \textbf{class and economic angle} to the caste angle.
\item \textbf{Daniel Thorner} (UP study): earlier we had Brahmin, Kshatriya, Vaishya, Shudra; \textbf{now we have Malik (owner), Kisan (cultivator) and Mazdoor (labourer)} --- the upper caste became Malik, the middle caste Kisan, the low caste Mazdoor --- \textbf{caste defined class; there was no change}.
\item \textbf{Kathleen Gough} (Kumbapettai, Tamil Nadu): \textbf{modernization of the agrarian structure did not help the poor castes --- it consolidated class divisions} (the village moved from caste-division to a strict class hierarchy).
\end{itemize}
\end{KeyConcept}

\begin{CriticalPoint}
\begin{itemize}
\item The Marxist scholars studied 20 villages in Punjab: \textbf{land reforms and the Green Revolution} (the two big post-independence programmes) at the surface seem \textbf{development-oriented} --- but the upper caste, with its social and cultural capital, bought the heavy machinery and high-yielding seeds, becoming 'super-upper-class', while the Dalit peasants fell further. The \textbf{rich farmers captured the seed, credit and irrigation cooperatives}.
\item The \textbf{failure of these policies} is what produced the \textbf{Maoist/Naxalite movement} --- Naxalism exists because land reform was not successful.
\end{itemize}
\end{CriticalPoint}

\section{Moffatt's Replication and the Contribution of Village Studies}

\begin{KeyConcept}
\begin{itemize}
\item \textbf{Moffatt} (recap) provided a study that is neither about integration nor conflict: the excluded untouchables, instead of conflicting, \textbf{replicate} the same caste structure that excluded them --- the principles of \textbf{replication and subordination}.
\item \textbf{Conclusion}: village studies not only analysed development, social organisation and the strategies of the excluded, but \textbf{contributed concepts to Indian sociology} --- \textbf{dominant caste, Sanskritization, dominant individual, replicability, consensus} --- which we still use (e.g. to explain reservation demands by the dominant caste).
\end{itemize}
\end{KeyConcept}

\begin{QuickRevision}
\begin{itemize}
\item 1950s: self-sufficiency doctrine questioned (Srinivas \& A. M. Shah, 'Myth of Self-Sufficiency'; Beteille's multi-language Brahmins; S. C. Dube's village-town interdependence; Phule's internal drain of wealth).
\item S. C. Dube (Shamirpet, 1955): multi-caste/multi-religious village; jajmani across religions; six criteria for a village typology; three broad features (unit; inter-ethnic dependency; villages integrated).
\item Village community organised by kin/caste (family $\rightarrow$ kin $\rightarrow$ lineage $\rightarrow$ relatives $\rightarrow$ sub-caste $\rightarrow$ varna) and by territory (individual $\rightarrow$ family $\rightarrow$ village $\rightarrow$ caste org $\rightarrow$ region $\rightarrow$ nation).
\item Beteille: Weberian model (symmetrical $\rightarrow$ asymmetrical) as an ideal type for comparison.
\item Marxist studies (1970s): Thorner (Malik/Kisan/Mazdoor --- no change), Kathleen Gough (Kumbapettai --- class consolidation); land reform + Green Revolution favoured rich farmers $\rightarrow$ Naxalism.
\item Moffatt: replication/subordination (not integration, not conflict); village studies gave concepts (dominant caste, Sanskritization, dominant individual, replicability, consensus).
\end{itemize}
\end{QuickRevision}

\section{Current Affairs  to  Sociology: The MTP Act and the Female Body}

\begin{ExampleBox}
\begin{itemize}
\item The \textbf{Medical Termination of Pregnancy (MTP) Act} --- a legislation to reduce illegal abortion and consequent maternal mortality/morbidity --- is a worked example of extracting sociology from a current-affairs topic.
\item The changes: the \textbf{gestational limit was raised from 20 to 24 weeks} (one doctor's approval up to 20 weeks, two up to 24); and '\textbf{pregnant wife}' was changed to '\textbf{pregnant woman}' (so unmarried women too can now terminate).
\item \textbf{Sociological reading}: (1) the gestational limit means the \textbf{state decides when a woman may keep or remove a child} --- the female body becomes a site of control --- the \textbf{medicalization of the female body}, supervised by (male) doctors. \textbf{Adrienne Rich} ('Of Woman Born') argues that this control of reproduction perpetuates patriarchy, with \textbf{male doctors replacing female midwives} (ANMs, ASHA workers).
\item (2) 'Pregnant \textbf{woman}' (not 'wife') \textbf{reinforces the gender binary / compulsory heterosexuality} --- excluding trans, intersex and non-binary people from safe abortion.
\item (3) After 24 weeks, a \textbf{medical board's approval} amounts to \textbf{bureaucratic surveillance} over women's bodies --- what \textbf{Michel Foucault} called \textbf{biopower} (control over the body).
\item (4) Treating '\textbf{fetal abnormality}' as a separate category reeks of \textbf{ableism} --- which legitimises the disabled as unproductive; but (as \textbf{Erving Goffman} says) 'normal' and 'abnormal' is only a matter of \textbf{perspective}.
\end{itemize}
\end{ExampleBox}

\begin{QuickRevision}
\begin{itemize}
\item MTP Act: gestational limit 20 $\rightarrow$ 24 weeks; 'pregnant wife' $\rightarrow$ 'pregnant woman'.
\item Medicalization of the female body (state/doctors control the body); Adrienne Rich ('Of Woman Born' --- male doctors replacing female midwives).
\item Compulsory heterosexuality / gender binary; Foucault's biopower (medical board surveillance); ableism (Goffman's normal/abnormal = perspective).
\end{itemize}
\end{QuickRevision}

